\section{Mo's Algorithm}
\whenuse{Muchas consultas offline sobre subarreglos, sin updates o con una
extensi\'on controlada, cuando agregar/quitar un extremo es barato.}

Reordena consultas est\'aticas de rango para mover dos punteros pocas veces.
Los rangos son semiabiertos $[l,r)$. El estado debe admitir agregar y quitar
un elemento en $O(1)$ o costo peque\~no.

\complexity{aproximadamente $O((n+q)\sqrt n)$ movimientos; memoria $O(n+q)$.}

\begin{lstlisting}[style=cpp]
struct MoQuery { int left, right, index; };

template<class Add, class Remove, class GetAnswer>
vector<ll> mos_algorithm(int n, vector<MoQuery> queries,
                         Add add, Remove remove, GetAnswer get_answer) {
    int block_size = max(1, (int)sqrt(max(1, n)));
    sort(queries.begin(), queries.end(), [&](const MoQuery &a, const MoQuery &b) {
        int block_a = a.left / block_size, block_b = b.left / block_size;
        if (block_a != block_b) return block_a < block_b;
        return block_a & 1 ? a.right > b.right : a.right < b.right;
    });

    vector<ll> answers(queries.size());
    int current_left = 0, current_right = 0;
    for (const MoQuery &query : queries) {
        while (current_left > query.left) add(--current_left);
        while (current_right < query.right) add(current_right++);
        while (current_left < query.left) remove(current_left++);
        while (current_right > query.right) remove(--current_right);
        answers[query.index] = get_answer();
    }
    return answers;
}
\end{lstlisting}

Ejemplos: cantidad de valores distintos, suma de frecuencias al cuadrado y
consultas offline sobre frecuencias. Para actualizaciones existe Mo con tiempo.

Ejemplo: cantidad de valores distintos en cada rango. \texttt{add(index)} agrega
\texttt{values[index]} a la ventana actual, \texttt{remove(index)} lo quita y
\texttt{get\_answer()} devuelve el estado mantenido.

\begin{lstlisting}[style=cpp]
vector<int> values={1,2,1,3,2};
vector<MoQuery> queries={{0,3,0},{2,5,1}}; // [0,3), [2,5)
unordered_map<int,int> frequency;
ll distinct=0;
auto add=[&](int index) {
    if(frequency[values[index]]++==0) ++distinct;
};
auto remove=[&](int index) {
    if(--frequency[values[index]]==0) --distinct;
};
auto get_answer=[&]() { return distinct; };
vector<ll> answer=mos_algorithm(values.size(),queries,
                                add,remove,get_answer);
cout<<answer[0]<<' '<<answer[1]<<'\n'; // 2 3
\end{lstlisting}

Las consultas se procesan reordenadas, pero \texttt{query.index} restaura el
orden original de las respuestas. ``Offline'' significa que todas las consultas
deben conocerse antes de comenzar; Mo no acelera una sola consulta ni admite
updates comunes en esta versi\'on.
